\chapter{Introduction}
Remote meetings have become a cornerstone of modern collaboration, yet they often suffer from inefficiencies, particularly in convergent decision-making processes such as prioritization, problem-solving, and consensus-building.
Research indicates that up to 67\% of meetings are considered unproductive by participants, with common issues including off-topic discussions, unequal participation, and a lack of structure [Doodle, 2023; Harvard Business Review, 2022]

Although traditional facilitation strategies are effective, their practical implementation is limited by the fact that not every team has permanent access to a trained facilitator. Professional training and external consulting, while valuable, are costly and not continuously available.

The potential of AI-driven facilitation to address these gaps lies in its ability to provide real-time, context-sensitive support. Unlike static tools (e.g., timers or agendas), proactive AI systems can analyze conversations, detect deviations from predefined goals, and intervene with constructive suggestions—thereby enhancing goal orientation and reducing cognitive load on human moderators. Early adopters of AI in meetings, such as Otter.ai (real-time transcription) and Microsoft Teams’ AI recaps, demonstrate the feasibility of automated support, but these tools lack proactive, adaptive moderation tailored to convergent processes [Microsoft, 2024; Otter.ai, 2023]

\section{Background and Motivation}
\subsection{Inefficiencies in remote meetings: Challenges in convergent decision-making processes}
Convergent decision-making in remote meetings faces unique obstacles:
\begin{itemize}
  \item Off-topic discussions: Studies show that 30–50\% of meeting time is spent on irrelevant tangents, derailing progress [Rogelberg et al., 2006]. \cite{chenAreWeTrack2025}
  \item Unequal participation: Dominant voices can overshadow quieter team members, leading to groupthink and suboptimal outcomes [Nunamaker et al., 1991].
  \item Lack of structure: Without clear facilitation, meetings often lack timeboxing, role assignment, or explicit decision criteria, reducing efficiency [Schwartz, 2016].
\end{itemize}
Traditional solutions—such as hiring professional facilitators or training team members—are cost-prohibitive or scalability-limited. For example, a certified facilitator may charge €100–300/hour [International Association of Facilitators, 2024], making continuous support impractical for most organizations. Furthermore, not every team has permanent access to trained personnel, and ad-hoc facilitation often fails to address systemic issues like dysfunctional team dynamics or blind spots in group perception [Edmondson, 1999].

\subsection{Potential of AI-driven facilitation to enhance goal orientation }

AI can bridge this gap by:
\begin{itemize}
  \item Automating real-time analysis: Natural Language Processing (NLP) can transcribe and semantically analyze discussions to identify deviations from the meeting’s goal [Bommasani et al., 2021].
  \item Providing adaptive interventions: Context-sensitive prompts (e.g., "Let’s return to the agenda: prioritizing the top 3 solutions") can refocus conversations without human intervention.
  \item Democratizing participation: AI can nudge underrepresented voices (e.g., "Alex, you haven’t shared your perspective yet—would you like to add anything?") or summarize key points to ensure alignment.
\end{itemize}
Pilot studies on AI moderation tools (e.g., Humley’s AI meeting assistant) report 20–30\% reductions in meeting duration and improved participant satisfaction [Humley, 2023]. However, these tools often lack proactivity—reacting to issues rather than preventing them—and customization for specific use cases like convergent decision-making.

\section{Research Problem and Gap}
\subsection{Limitations of existing moderation tools in remote settings}
Current tools fall short in three key areas:
\begin{itemize}
  \item Passive functionality: Most AI tools (e.g., Zoom AI Companion, Google Meet’s recaps) focus on post-meeting summaries or transcription, not real-time moderation [Zoom, 2024; Google, 2023].
  \item Generic interventions: Solutions like Miro’s AI facilitator provide template-based guidance but fail to adapt to the meeting’s context or user preferences [Miro, 2024].
  \item Lack of empirical validation: Few studies evaluate user acceptance of proactive AI moderation, particularly in high-stakes decision-making scenarios [Benbasat and Barki, 2007].
\end{itemize}
\subsection{Need for context-sensitive, proactive AI systems }
To address these gaps, this thesis proposes a proactive AI moderation system that:
\begin{itemize}
  \item Detects thematic deviations in real time using NLP and meeting goal tracking.
  \item Intervenes constructively with personalized, non-intrusive prompts (e.g., private messages to individuals or group reminders).
  \item Validates design choices through user studies focusing on usability (SUS) and acceptance (TAM).
\end{itemize}
Such a system could reduce meeting fatigue, improve decision quality, and lower the barrier to effective facilitation—but its success hinges on user trust and perceived usefulness, which remain underexplored in literature.
\section{Research Objectives and Questions}
This thesis addresses the inefficiencies in goal-oriented remote meetings through a structured approach that integrates three complementary research objectives. The first objective (RQ1) investigates facilitating strategies and challenges in remote meetings to establish a foundation based on real-world practices and pain points. The second objective (RQ2) focuses on designing an AI-based moderation system that translates these insights into technical requirements, ensuring the system is context-aware and goal-oriented. The third objective (RQ3) evaluates user acceptance among knowledge workers.

The work bridges the gap between theoretical facilitation techniques and practical AI-driven solutions, ensuring that the system is both technically robust and user-centric.
\subsection{RQ1: Facilitating strategies and challenges in goal-oriented remote meetings}
\begin{quote}
  \textit{Which facilitating strategies are suitable for the goal-oriented management of remote meetings with convergent objectives (such as decision-making, prioritization, problem-solving)? What specific challenges typically arise in this context?}
\end{quote}
The first research question aims to identify effective strategies for managing convergent remote meetings, such as decision-making, prioritization, and problem-solving, while also analyzing the specific challenges that arise. To achieve this, the thesis conducts expert interviews with certified facilitators, agile coaches, and team leads to gather empirical insights. It also analyzes literature on facilitation techniques, including timeboxing and role assignment, as well as the dynamics of remote collaboration. The findings are compared with existing research to derive best practices and identify pain points, such as off-topic discussions and unequal participation, which AI could potentially address.
\subsection{RQ2: Design requirements for an AI-based moderation system}
\begin{quote}
  \textit{How should an AI-based moderation system be designed that autonomously detects thematic deviations based on a predefined meeting goal and real-time transcript analysis and promotes goal orientation through constructive, context-sensitive interventions?}
\end{quote}
The second research question focuses on developing a proactive AI system that detects thematic deviations in real time and intervenes contextually to maintain goal orientation. This objective is addressed by defining functional requirements, such as real-time transcript analysis and adaptive interventions, and non-functional requirements, including privacy, latency, and scalability. A system architecture is then designed with components for transcription, NLP-based analysis, and intervention logic, incorporating mechanisms for private nudges or group reminders. A prototype is implemented using NLP libraries like Hugging Face Transformers, and its technical feasibility is evaluated.
\subsection{RQ3: User acceptance among knowledge workers}
\begin{quote}
  \textit{How do knowledge workers in remote teams evaluate the usefulness, usability, and psychological effects of such an AI moderation tool, and what factors influence its acceptance?}
\end{quote}
The third research question assesses how knowledge workers perceive the usefulness, usability, and psychological impact of the AI moderation system, as well as the factors influencing its adoption. A usability study is conducted with remote teams, using frameworks such as the System Usability Scale and the Technology Acceptance Model. Quantitative data, including SUS scores and efficiency metrics, are collected alongside qualitative feedback from interviews that explore perceptions of control, trust, and pressure. The analysis also examines psychological effects, such as relief versus intrusion, and adoption barriers, including resistance to automation and privacy concerns.